\originalchapter{Quiz 1: 题目与解答}
试卷日期为2012年10月18日, 四题总分100. 本章依据 Alan Guth 的试卷与2012年10月21日发布的\textbf{官方解答中文整理}, 补充计算步骤并核对原件中的排版错误; 新增解释不冒充官方逐字原文. 原试卷附公式表, 仓库保留原表; 本稿将有关公式放在前述理论章节中.

\originalsection{短题}
\begin{exercise}
题1 (30分). (a, 10分) 用指标记号求 $\nabla\times(s\vct A)$, $s$ 为标量场. (b, 10分) 判断 $\ee=x\hat{\vct x}-y\hat{\vct y}$、$\ee=y\hat{\vct x}+x\hat{\vct y}$、$\ee=y\hat{\vct x}-x\hat{\vct y}$ 能否为静电场, 简述理由. (c, 10分) 整个 $xz$ 与 $yz$ 平面是导体, 求位于 $(x_0,y_0,0)$ 的电荷 $q$ 所受力, 电荷不位于导体面上.
\end{exercise}
\begin{solution}
(a) $\epsilon_{ijk}\partial_j(sA_k)=s\epsilon_{ijk}\partial_jA_k+\epsilon_{ijk}(\partial_js)A_k$, 故结果为 $s\nabla\times\vct A+\nabla s\times\vct A$.
(b) 三场的旋度依次为零、零、$-2\hat{\vct z}$. 前两个可分别由 $V=(y^2-x^2)/2$ 和 $V=-xy$ 产生, 第三个不满足静电旋度条件. 这里讨论局部可实现性, 不要求这些随位置增长的场是局域有限能电荷产生的全空间场.
(c) 两相交导体共有同一电势, 以它为零. 三镜像分别为 $-q$ 在 $(-x_0,y_0,0)$、$-q$ 在 $(x_0,-y_0,0)$、$+q$ 在 $(-x_0,-y_0,0)$. 两平面上成对抵消. 一般非零 $x_0,y_0$ 的力为
\[
\vct F=\frac{q^2}{4\pi\eps}\left[-\frac{x_0\hat{\vct x}}{4|x_0|^3}
-\frac{y_0\hat{\vct y}}{4|y_0|^3}
+\frac{x_0\hat{\vct x}+y_0\hat{\vct y}}{4(x_0^2+y_0^2)^{3/2}}\right].
\]
在第一象限化为官方解答的 $-\hat{\vct x}/(4x_0^2)$、$-\hat{\vct y}/(4y_0^2)$ 两项. 求真实电荷受力时不乘势能公式中的一半.
\end{solution}

\originalsection{柱形带电体}
\begin{exercise}
题2 (20分). 无限长实心圆柱半径 $a$、均匀体密度 $\rho$, 外面有半径 $b>a$ 的同轴薄圆柱面, 单位长度总电荷等大反向. (a, 7分) 求各处场; (b, 6分) 以外柱面为零势求势; (c, 7分) 求单位长度场能.
\end{exercise}
\begin{solution}
高斯柱面的净通量为 $2\pi s\ell E_s$. 内柱包围电荷 $\pi s^2\ell\rho$, 柱外到 $b$ 之间为 $\pi a^2\ell\rho$, $b$ 外净电荷为零, 所以
\[
\ee=\frac\rho{2\eps}\begin{cases}s\hat{\vct s},&s<a,\\ a^2\hat{\vct s}/s,&a<s<b,\\0,&s>b.\end{cases}
\]
官方解答的个别分段标签误排成 $s>a$; 中文式以原物理区域 $s>b$ 为准. 由 $V(s)=\int_s^bE_s\dd s$,
\[
V(s)=\begin{cases}
\rho a^2\ln(b/a)/(2\eps)+\rho(a^2-s^2)/(4\eps),&s<a,\\
\rho a^2\ln(b/s)/(2\eps),&a<s<b,\\0,&s>b.
\end{cases}
\]
单位长度能量为 $\pi\eps\int_0^bE_s^2s\dd s$. 内部积分给 $\pi\rho^2a^4/(16\eps)$, 外环给 $\pi\rho^2a^4\ln(b/a)/(4\eps)$, 因而
$U/\ell=\pi\rho^2a^4[\frac14+\ln(b/a)]/(4\eps)$. 内部非零能量与金属内柱电容器不同, 因为这里内柱是均匀体电荷而非等势导体.
\end{solution}

\originalsection{带电短线的多极}
\begin{exercise}
题3 (20分). 短线位于 $z$ 轴、中心在原点, 总电荷 $Q\ne0$, 线密度为偶函数. 给定 $\ell_0^2=Q^{-1}\int z^2\lambda(z)\dd z$. (a, 10分) 求电偶极和四极矩; (b, 10分) 写远势到四极阶.
\end{exercise}
\begin{solution}
偶函数使一阶矩为零, 线在 $z$ 轴使横向坐标为零. 因而 $\pp=0$, $Q_{xx}=Q_{yy}=-Q\ell_0^2$, $Q_{zz}=2Q\ell_0^2$, 非对角项为零. 与 $n_in_j$ 收缩得到 $Q\ell_0^2(3\cos^2\theta-1)$, 所以
\[
V(r,\theta)=\frac Q{4\pi\eps}\left[\frac1r+
\frac{\ell_0^2}{2r^3}(3\cos^2\theta-1)\right]+O(r^{-5}).
\]
偶对称同时消掉奇数阶, 因而余项从四阶势 $r^{-5}$ 开始; 该估计假设源局域且相应矩有限.
\end{solution}

\originalsection{非均匀球壳}
\begin{exercise}
题4 (30分). 半径 $R$ 的球壳中心在原点, 无穷远零势. 依次求内外电势: (a, 10分) 表面势为 $V_0\sin\theta\cos\phi$; (b, 10分) 表面势为 $V_0\sin^2\theta\sin^2\phi$; (c, 10分) 改为给定表面电荷 $\sigma_0\sin^2\theta\sin^2\phi$.
\end{exercise}
\begin{solution}
(a) $\sin\theta\cos\phi=n_x$ 为一阶谐函数, 内势 $V_0(r/R)n_x$, 外势 $V_0(R/r)^2n_x$.
(b) 令 $h=n_y^2-1/3$. $n_y^2=1/3+h$ 分成零阶和二阶, 故
\[
V_\mathrm{in}=V_0[1/3+(r/R)^2h],\quad
V_\mathrm{out}=V_0[R/(3r)+(R/r)^3h].
\]
(c) 若表面电荷某阶为 $\sigma_\ell F_\ell$, 共同表面势为 $A_\ell F_\ell$, 内外径向导数差给
$\sigma_\ell=\eps(2\ell+1)A_\ell/R$. 因此
\[
V_\mathrm{in}=\frac{\sigma_0R}\eps[1/3+(r/R)^2h/5],\quad
V_\mathrm{out}=\frac{\sigma_0R}\eps[R/(3r)+(R/r)^3h/5].
\]
代入 $r=R$ 检查势连续, 再核对法向跳跃恢复 $\sigma_0n_y^2$. 不能把(b)的势直接乘常数当(c), 因为零阶与二阶的电荷到电势系数不同.
\end{solution}
